\documentclass[10pt,a4paper]{amsart}
\usepackage{amsmath,amssymb,amsthm,graphicx}
% Retained preamble for the Stacks Project chapter {CH} (shared project preamble).
% The chapter file pinned-original-source.tex only carries \input{{preamble}} on its first line,
% so no preamble text lives inside the chapter itself.  The Stacks Project uses one single shared
% preamble for every chapter: the snapshot stages it three times (sources/stacks-chow/preamble.tex,
% sources/stacks-etale-cohomology/preamble.tex and focused-seeds/AGP/sources/stacks-snapshot-20261011/
% preamble.tex) and those copies are identical in size.  The block below is copied mechanically,
% line for line, from sources/stacks-chow/preamble.tex lines 14-18 and 154-217: the Xy-pic setup and
% the native theorem declarations together with the mathematical macros.  Nothing was added, removed
% or redefined; the wrapper class and the amsmath/amssymb/amsthm/graphicx base packages are supplied
% by the assembler.
% For commutative diagrams we use Xy-pic
% Chapter: stacks-cohomology.tex -- Cohomology of Sheaves.
\usepackage[all]{xy}

% We use 2cell for 2-commutative diagrams.
\xyoption{2cell}

% Theorem environments.
%
\theoremstyle{plain}
\newtheorem{theorem}[subsection]{Theorem}
\newtheorem{proposition}[subsection]{Proposition}
\newtheorem{lemma}[subsection]{Lemma}

\theoremstyle{definition}
\newtheorem{definition}[subsection]{Definition}
\newtheorem{example}[subsection]{Example}
\newtheorem{exercise}[subsection]{Exercise}
\newtheorem{situation}[subsection]{Situation}

\theoremstyle{remark}
\newtheorem{remark}[subsection]{Remark}
\newtheorem{remarks}[subsection]{Remarks}

\numberwithin{equation}{subsection}

% Macros
%
\def\lim{\mathop{\mathrm{lim}}\nolimits}
\def\colim{\mathop{\mathrm{colim}}\nolimits}
\def\Spec{\mathop{\mathrm{Spec}}}
\def\Hom{\mathop{\mathrm{Hom}}\nolimits}
\def\Ext{\mathop{\mathrm{Ext}}\nolimits}
\def\SheafHom{\mathop{\mathcal{H}\!\mathit{om}}\nolimits}
\def\SheafExt{\mathop{\mathcal{E}\!\mathit{xt}}\nolimits}
\def\Sch{\mathit{Sch}}
\def\Mor{\mathop{\mathrm{Mor}}\nolimits}
\def\Ob{\mathop{\mathrm{Ob}}\nolimits}
\def\Sh{\mathop{\mathit{Sh}}\nolimits}
\def\NL{\mathop{N\!L}\nolimits}
\def\CH{\mathop{\mathrm{CH}}\nolimits}
\def\proetale{{pro\text{-}\acute{e}tale}}
\def\etale{{\acute{e}tale}}
\def\QCoh{\mathit{QCoh}}
\def\Ker{\mathop{\mathrm{Ker}}}
\def\Im{\mathop{\mathrm{Im}}}
\def\Coker{\mathop{\mathrm{Coker}}}
\def\Coim{\mathop{\mathrm{Coim}}}

% Boxtimes
%
\DeclareMathSymbol{\boxtimes}{\mathbin}{AMSa}{"02}

%
% Macros for moduli stacks/spaces
%
\def\QCohstack{\mathcal{QC}\!\mathit{oh}}
\def\Cohstack{\mathcal{C}\!\mathit{oh}}
\def\Spacesstack{\mathcal{S}\!\mathit{paces}}
\def\Quotfunctor{\mathrm{Quot}}
\def\Hilbfunctor{\mathrm{Hilb}}
\def\Curvesstack{\mathcal{C}\!\mathit{urves}}
\def\Polarizedstack{\mathcal{P}\!\mathit{olarized}}
\def\Complexesstack{\mathcal{C}\!\mathit{omplexes}}
% \Pic is the operator that assigns to X its picard group, usage \Pic(X)
% \Picardstack_{X/B} denotes the Picard stack of X over B
% \Picardfunctor_{X/B} denotes the Picard functor of X over B
\def\Pic{\mathop{\mathrm{Pic}}\nolimits}
\def\Picardstack{\mathcal{P}\!\mathit{ic}}
\def\Picardfunctor{\mathrm{Pic}}
\def\Deformationcategory{\mathcal{D}\!\mathit{ef}}
\usepackage{cleveref}
\crefname{originalreference}{original reference}{original references}
\Crefname{originalreference}{Original reference}{Original references}
\makeatletter
\newlabel{modules-lemma-finite-presentation-stalk-free}{{0B8J}{source}{}{sourcepointer-1}{}}
\newlabel{modules-lemma-finite-presentation-stalk-free@cref}{{[originalreference][0][]0B8J}{[1][0][]source}{}{}{}}
\makeatother
\crefname{originallemma}{original lemma}{original lemmas}
\Crefname{originallemma}{Original lemma}{Original lemmas}
\makeatletter
\newlabel{lemma-local-lift-map}{{08C6}{source}{}{sourcepointer-2}{}}
\newlabel{lemma-local-lift-map@cref}{{[originallemma][0][]08C6}{[1][0][]source}{}{}{}}
\makeatother
\title{The locus where all cohomology sheaves of a perfect complex are finite free is open and is the maximal such open subset}
\author{The Stacks Project Authors}
\date{stacks-project:cohomology.tex:local-snapshot-20261011; GNU FDL 1.2 or later}
\begin{document}
\maketitle
\noindent\textit{Editorial note.} All delivered mathematical text is the author's own, copied verbatim from the pinned Stacks Project chapter file cohomology.tex; no statement or proof step was rewritten, reordered or removed.\par\medskip
\noindent\textit{Editorial note.} The counted span cites two cross-chapter labels, each recorded in assembly.external\_references with its canonical tag as a navigation pointer.\par\medskip
\noindent\textit{Editorial note.} The proof cites no bibliography item inside the counted span.\par\medskip
\noindent\textit{Allowed standard prerequisites (editorial).}\begin{itemize}
\item A finitely presented module with free stalk is free on a neighbourhood. Stacks Project, Modules on Sites, Lemma modules-lemma-finite-presentation-stalk-free (tag 0B8J). Spreads freeness of the stalk of the top cohomology sheaf to an open neighbourhood.
\item Lifting a direct sum decomposition of a quotient. Stacks Project, Cohomology of Sheaves, Lemma lemma-local-lift-map (tag 08C6). Produces the direct sum decompositions of the strictly perfect complex used in the induction step.
\end{itemize}\medskip
\ifdefined\hypertarget\hypertarget{sourcepointer-1}{}\fi
\noindent\textit{Source reference.} Original reference 0B8J: sources/stacks-chow/tags.txt, source lines 14142--14142. Cross-chapter standard prerequisite cited by the author inside the counted span; it is used as a pointer and its own proof is not part of this item.\par\medskip
\ifdefined\hypertarget\hypertarget{sourcepointer-2}{}\fi
\noindent\textit{Source reference.} Original lemma 08C6: sources/stacks-chow/tags.txt, source lines 10591--10591. Cross-chapter standard prerequisite cited by the author inside the counted span; it is used as a pointer and its own proof is not part of this item.\par\medskip
\begin{lemma}
\label{lemma-perfect-max-open-coh-loc-free}
Let $(X, \mathcal{O}_X)$ be a ringed space. Let $E$ in $D(\mathcal{O}_X)$
be perfect. Assume that all stalks $\mathcal{O}_{X, x}$ are local rings.
Then the set
$$
U =
\{x \in X \mid
H^i(E)_x\text{ is a finite free }
\mathcal{O}_{X, x}\text{-module for all }i\in \mathbf{Z}\}
$$
is open in $X$ and is the maximal open set $U \subset X$ such that
$H^i(E)|_U$ is finite locally free for all $i \in \mathbf{Z}$.
\end{lemma}

\begin{proof}
Note that if $V \subset X$ is some open such that $H^i(E)|_V$
is finite locally free for all $i \in \mathbf{Z}$ then $V \subset U$. 
Let $x \in U$. We will show that an open neighbourhood of $x$
is contained in $U$ and that $H^i(E)$ is finite locally free
on this neighbourhood for all $i$. This will finish the proof.
During the proof we may (finitely many times) replace
$X$ by an open neighbourhood of $x$.
Hence we may assume $E$ is represented
by a strictly perfect complex $\mathcal{E}^\bullet$. Say
$\mathcal{E}^i = 0$ for $i \not \in [a, b]$. We will
prove the result by induction on $b - a$. The module
$H^b(E) = \Coker(d^{b - 1} : \mathcal{E}^{b - 1} \to \mathcal{E}^b)$
is of finite presentation. Since $H^b(E)_x$ is finite free,
we conclude $H^b(E)$ is finite free in an open neighbourhood of $x$ by
Modules, Lemma \ref{modules-lemma-finite-presentation-stalk-free}.
Thus after replacing $X$ by a (possibly smaller) open
neighbourhood we may assume we have a direct sum decomposition
$\mathcal{E}^b = \Im(d^{b - 1}) \oplus H^b(E)$ and
$H^b(E)$ is finite free, see Lemma \ref{lemma-local-lift-map}.
Doing the same argument again, we see that we may
assume $\mathcal{E}^{b - 1} = \Ker(d^{b - 1}) \oplus \Im(d^{b - 1})$.
The complex
$\mathcal{E}^a \to \ldots \to \mathcal{E}^{b - 2} \to \Ker(d^{b - 1})$
is a strictly perfect complex representing a perfect
object $E'$ with $H^i(E) = H^i(E')$ for $i \not = b$.
Hence we conclude by our induction hypothesis.
\end{proof}

\end{document}
